\SixthArticle{41}

Let us next set $P=x^\mu P'$, so that
\[
\begin{aligned}
\frac{\mathrm{d}P}{\mathrm{d}x}&=\mu x^{\mu-1}P'+x^\mu\frac{\mathrm{d}P'}{\mathrm{d}x}\\
\frac{\mathrm{dd}P}{\mathrm{d}x^2}&=(\mu\mu-\mu)x^{\mu-2}P'+2\mu x^{\mu-1}\frac{\mathrm{d}P'}{\mathrm{d}x}+x^\mu\frac{\mathrm{dd}P'}{\mathrm{d}x^2}
\end{aligned}
\]
\SourceLeaf{219}
Substituting these values in 80 and dividing by $x^{\mu-1}$ gives
\[
\begin{aligned}
0={}&P'\{\alpha\beta x-(\gamma-(\alpha+\beta+1)x)\mu-(1-x)(\mu\mu-\mu)\}\\
&-\frac{\mathrm{d}P'}{\mathrm{d}x}\{\gamma-(\alpha+\beta+1)x+2\mu(1-x)\}x\\
&-\frac{\mathrm{dd}P'}{\mathrm{d}x^2}(xx-x^3)
\end{aligned}
\]
The multiplier of $P'$ in this formula becomes divisible by $x$ on setting $\mu=0$ or $\mu=1-\gamma$; the latter value gives
\[
\begin{aligned}
0={}&P'(\alpha\beta+\alpha+\beta+1-2\gamma-\alpha\gamma-\beta\gamma+\gamma\gamma)\\
&-\frac{\mathrm{d}P'}{\mathrm{d}x}(2-\gamma-(\alpha+\beta+3-2\gamma)x)\\
&-\frac{\mathrm{dd}P'}{\mathrm{d}x^2}(x-xx)
\end{aligned}
\]
Comparing this equation with 80, whose form is exactly similar, shows that what were $P$, $\alpha$, $\beta$, $\gamma$ there are $P'$, $\alpha+1-\gamma$, $\beta+1-\gamma$, $2-\gamma$ here. Since we have given the general solution of that equation, it is evident that $P'$ will have the form
\[
\begin{aligned}
P'={}&MF(\alpha+1-\gamma,\beta+1-\gamma,2-\gamma,x)\\
&+NF(\alpha+1-\gamma,\beta+1-\gamma,\alpha+\beta+1-\gamma,1-x)
\end{aligned}
\]
where $M$, $N$ denote constant quantities, or
\[
\begin{aligned}
&\text{[83]}\qquad F(\alpha,\beta,\gamma,x)\\&=Mx^{1-\gamma}F(\alpha+1-\gamma,\beta+1-\gamma,2-\gamma,x)\\
&+Nx^{1-\gamma}F(\alpha+1-\gamma,\beta+1-\gamma,\alpha+\beta+1-\gamma,1-x)
\end{aligned}
\]
where the constants $M$, $N$ will depend on the elements $\alpha$, $\beta$, $\gamma$.

\SixthArticle{42}

From equation 82 follows
\[
\begin{aligned}
&F(\alpha+1-\gamma,\beta+1-\gamma,2-\gamma,x)\\&=(1-x)^{\gamma-\alpha-\beta}F(1-\alpha,1-\beta,2-\gamma,x)\\[5pt]
&F(\alpha+1-\gamma,\beta+1-\gamma,\alpha+\beta+1-\gamma,1-x)\\&=x^{\gamma-1}F(\alpha,\beta,\alpha+\beta+1-\gamma,1-x)
\end{aligned}
\]
\Needspace{16\baselineskip}
Thus, setting
\[
\frac{1}{N}=f(\alpha,\beta,\gamma),\qquad-\frac{M}{N}=g(\alpha,\beta,\gamma)
\]
\SourceLeaf{220}
equation 83 becomes
\[
\begin{aligned}
&F(\alpha,\beta,\alpha+\beta+1-\gamma,1-x)\\&=f(\alpha,\beta,\gamma)F(\alpha,\beta,\gamma,x)\\
&\quad+g(\alpha,\beta,\gamma)(1-x)^{\gamma-\alpha-\beta}x^{1-\gamma}F(1-\alpha,1-\beta,2-\gamma,x)
\end{aligned}
\]
With the aid of formula 82, the same equation can also be given the form
\[
\begin{aligned}
&x^{1-\gamma}F(\alpha+1-\gamma,\beta+1-\gamma,\alpha+\beta+1-\gamma,1-x)=\\
&f(\alpha,\beta.\gamma)(1-x)^{\gamma-\alpha-\beta}F(\gamma-\alpha,\gamma-\beta,\gamma,x)
\\&+g(\alpha,\beta,\gamma)x^{1-\gamma}F(\alpha+1-\gamma,\beta+1-\gamma,2-\gamma,x)
\end{aligned}
\]
or, dividing by \(x^{1-\gamma}\) and replacing \(\alpha,\beta,\gamma\) respectively by \(\alpha+1-\gamma,\beta+1-\gamma,2-\gamma\),
\[
\begin{aligned}
&F(\alpha,\beta,\alpha+\beta+1-\gamma,1-x)\\
&\quad=g(\alpha+1-\gamma,\beta+1-\gamma,2-\gamma)F(\alpha,\beta,\gamma,x)\\
&\qquad+f(\alpha+1-\gamma,\beta+1-\gamma,2-\gamma)\\&\qquad(1-x)^{\gamma-\alpha-\beta}x^{1-\gamma}F(1-\alpha,1-\beta,2-\gamma,x)
\end{aligned}
\]
Since this must be identical with the preceding formula, we have
\[
g(\alpha,\beta,\gamma)=f(\alpha+1-\gamma,\beta+1-\gamma,2-\gamma)
\]
and therefore
\[
\begin{aligned}
\text{[84]}\quad &F(\alpha,\beta,\alpha+\beta+1-\gamma,1-x)\\
&=f(\alpha,\beta,\gamma)F(\alpha,\beta,\gamma,x)\\
&\quad+f(\alpha+1-\gamma,\beta+1-\gamma,2-\gamma)\\&\qquad(1-x)^{\gamma-\alpha-\beta}x^{1-\gamma}F(1-\alpha,1-\beta,2-\gamma,x)
\end{aligned}
\]
\SixthArticle{43}
Now, to determine the nature of the function \(f(\alpha,\beta,\gamma)\), let us set \(x=0\). It is then evident that \(F(\alpha,\beta,\gamma,x)=1.\ x^{1-\gamma}=0\) whenever \(1-\gamma\) is a positive quantity. But by equation 48 we have
\[
F(\alpha,\beta,\alpha+\beta+1-\gamma,1)=\frac{\Pi(\alpha+\beta-\gamma)\Pi(-\gamma)}{\Pi(\alpha-\gamma)\Pi(\beta-\gamma)}
\]
Thus, under the same restriction, it has been proved that
\[
\text{[85]}\qquad f(\alpha,\beta,\gamma)=\frac{\Pi(\alpha+\beta-\gamma)\Pi(-\gamma)}{\Pi(\alpha-\gamma)\Pi(\beta-\gamma)}
\]
\Needspace{19\baselineskip}
We prove that this formula is general as follows.
\par\SourceLeaf{221}
Differentiating equation 84 gives
\[
\begin{aligned}
&-\frac{\alpha\beta}{\alpha+\beta+1-\gamma}F(\alpha+1,\beta+1,\alpha+\beta+2-\gamma,1-x)\\
&=\frac{\alpha\beta}{\gamma}f(\alpha,\beta,\gamma)F(\alpha+1,\beta+1,\gamma+1.\ x)\\
&\quad+f(\alpha+1-\gamma,\beta+1-\gamma,2-\gamma)(1-x)^{\gamma-\alpha-\beta-1}x^{-\gamma}\\&\qquad\Bigl\{((1-\gamma)(1-x)-(\gamma-\alpha-\beta)x)F(1-\alpha,1-\beta,2-\gamma,x)\\
&\qquad+\frac{(1-\alpha)(1-\beta)}{2-\gamma}(x-xx)F(2-\alpha,2-\beta,3-\gamma,x)\Bigr\}
\end{aligned}
\]
But by formula IX of Article 10, replacing \(\alpha,\beta,\gamma\) by \(-\alpha,-\beta,1-\gamma\) gives
\[
\begin{aligned}
&(1-\gamma)(2-\gamma)F(-\alpha,-\beta,1-\gamma,x)\\&=(2-\gamma)(1-\gamma+(\alpha+\beta-1)x)F(1-\alpha,1-\beta,2-\gamma,x)\\
&\quad+(1-\alpha)(1-\beta)(x-xx)F(2-\alpha,2-\beta,3-\gamma,x)
\end{aligned}
\]
and the preceding equation consequently becomes
\[
\begin{aligned}
&F(\alpha+1,\beta+1,\alpha+\beta+2-\gamma,1-x)\\
&=-\frac{\alpha+\beta+1-\gamma}{\gamma}f(\alpha,\beta,\gamma)F(\alpha+1,\beta+1,\gamma+1.\ x)\\
&\quad-\frac{(\alpha+\beta+1-\gamma)(1-\gamma)}{\alpha\beta}f(\alpha+1-\gamma,\beta+1-\gamma,2-\gamma)\\&\qquad(1-x)^{\gamma-\alpha-\beta-1}x^{-\gamma}F(-\alpha,-\beta,1-\gamma,x)
\end{aligned}
\]
On the other hand, in equation 84, replacing \(\alpha,\beta,\gamma\) by \(\alpha+1,\beta+1,\gamma+1\) gives
\[
\begin{aligned}
&F(\alpha+1,\beta+1,\alpha+\beta+2-\gamma,1-x)\\
&\quad=f(\alpha+1,\beta+1,\gamma+1)F(\alpha+1,\beta+1,\gamma+1,x)\\
&\qquad+f(\alpha+1-\gamma,\beta+1-\gamma,1-\gamma)\\&\qquad(1-x)^{\gamma-\alpha-\beta-1}x^{-\gamma}F(-\alpha,-\beta,1-\gamma,x)
\end{aligned}
\]
Since it is readily seen that these two equations must be identical, we have in general
\[
f(\alpha+1,\beta+1,\gamma+1)=\frac{\alpha+\beta+1-\gamma}{-\gamma}f(\alpha,\beta,\gamma)
\]
or, replacing \(\alpha,\beta,\gamma\) by \(\alpha-1,\beta-1,\gamma-1\),
\[
\begin{aligned}
f(\alpha,\beta,\gamma)&=\frac{\alpha+\beta-\gamma}{1-\gamma}f(\alpha-1,\beta-1,\gamma-1)\\
&=\frac{\alpha+\beta-\gamma\,.\,\alpha+\beta-\gamma-1}{1-\gamma\,.\,2-\gamma}f(\alpha-2,\beta-2,\gamma-2)
\end{aligned}
\]
etc. From this it is easy to conclude that, in general, for any integer value of \(k\),
\[
f(\alpha,\beta,\gamma)=\frac{\Pi(\alpha+\beta-\gamma)\Pi(-\gamma)}{\Pi(\alpha+\beta-\gamma-k)\Pi(k-\gamma)}.f(\alpha-k,\beta-k,\gamma-k)
\]
\par\SourceLeaf{222}
But whenever $1-(\gamma-k)$, or $k+1-\gamma$, is a positive quantity, we have proved (formula 85) that
\[
f(\alpha-k,\beta-k,\gamma-k)
=\frac{\Pi(\alpha+\beta-\gamma-k)\Pi(k-\gamma)}{\Pi(\alpha-\gamma)\Pi(\beta-\gamma)}
\]
Thus, since $k$ can always be taken large enough, whatever $\gamma$ may be, to make $k+1-\gamma$ a positive quantity, we have \textit{in general}
\[
f(\alpha,\beta,\gamma)
=\frac{\Pi(\alpha+\beta-\gamma)\Pi(-\gamma)}{\Pi(\alpha-\gamma)\Pi(\beta-\gamma)}
\]
and hence
\[
f(\alpha+1-\gamma,\beta+1-\gamma,2-\gamma)
=\frac{\Pi(\alpha+\beta-\gamma)\Pi(\gamma-2)}{\Pi(\alpha-1)\Pi(\beta-1)}
\]
so that our formula becomes
\[
\begin{aligned}[c]
[86]\quad &F(\alpha,\beta,\alpha+\beta+1-\gamma,1-x)\\
&=\frac{\Pi(\alpha+\beta-\gamma)\Pi(-\gamma)}{\Pi(\alpha-\gamma)\Pi(\beta-\gamma)}F(\alpha,\beta,\gamma,x)\\
&\quad+\frac{\Pi(\alpha+\beta-\gamma)\Pi(\gamma-2)}{\Pi(\alpha-1)\Pi(\beta-1)}x^{1-\gamma}\\&\qquad(1-x)^{\gamma-\alpha-\beta}F(1-\alpha,1-\beta,2-\gamma,x)
\end{aligned}
\]
or, replacing $\gamma$ by $\alpha+\beta+1-\gamma$,
\[
\begin{aligned}[c]
[87]\quad &F(\alpha,\beta,\gamma,1-x)\\
&=\frac{\Pi(\gamma-1)\Pi(\gamma-\alpha-\beta-1)}{\Pi(\gamma-\alpha-1)\Pi(\gamma-\beta-1)}F(\alpha,\beta,\alpha+\beta+1-\gamma,x)\\
&\quad+\frac{\Pi(\gamma-1)\Pi(\alpha+\beta-\gamma-1)}{\Pi(\alpha-1)\Pi(\beta-1)}x^{\gamma-\alpha-\beta}\\&\qquad(1-x)^{1-\gamma}F(1-\alpha,1-\beta,\gamma+1-\alpha-\beta,x)
\end{aligned}
\]

If preferred, one may write

in formula 86
\[
\begin{aligned}&\text{for }(1-x)^{\gamma-\alpha-\beta}F(1-\alpha,1-\beta,2-\gamma,x)\\&\ .\ .\ .\ .\ F(\alpha+1-\gamma,\beta+1-\gamma,2-\gamma,x)\end{aligned}
\]
in formula 87
\[
\begin{aligned}&\text{for }(1-x)^{1-\gamma}F(1-\alpha,1-\beta,\gamma+1-\alpha-\beta,x)\\&\ .\ .\ .\ .\ F(\gamma-\alpha,\gamma-\beta,\gamma+1-\alpha-\beta,x)\end{aligned}
\]

\SixthArticle{44}

Thus, whenever the fourth element in a series of our form is assigned a value between $0{,}5$ and $1$, its slower convergence can be remedied by the preceding formulas: they split it into two other similar series,
\SourceLeaf{223}
which converge the more rapidly, the more slowly the original series converged. But the special cases in which this transformation fails must be excepted, namely whenever, in the series to be transformed, the difference between the third element and the sum of the first two elements is an integer. For if, in formula 86, $\gamma$ $=0$ or is equal to a negative integer, $F(\alpha,\beta,\gamma,x)$ evidently becomes an unusable series (Article 2), and the factor $\Pi(\gamma-2)$ infinite. If, however, $\gamma$ is a positive integer greater than unity, $F(1-\alpha,1-\beta,2-\gamma,x)$ and $F(\alpha+1-\gamma,\beta+1-\gamma,2-\gamma,x)$ become unusable series, and $\Pi(-\gamma)$ infinite. Finally, if $\gamma=1$, the two transformed series $F(\alpha,\beta,\gamma,x)$ and $F(1-\alpha,1-\beta,2-\gamma,x)$, or $F(\alpha+1-\gamma,\beta+1-\gamma,2-\gamma,x)$, which therefore becomes identical with $F(\alpha,\beta,\gamma,x)$, do not suffer from this difficulty; nevertheless, the transformation is of no use, since each of the two transformed series is multiplied by the infinite coefficient $\Pi(-1)$. It will therefore be worthwhile to show how, in these cases too, slower convergence can be changed into more rapid convergence.

